\documentclass[11pt,a4paper,numbers=noenddot]{scrartcl}

\usepackage[T1]{fontenc}
\usepackage[utf8]{inputenc}
\usepackage[nochapters,eulerchapternumbers,beramono,pdfspacing,dottedtoc]{classicthesis}
\usepackage{amsmath,amssymb,bm}
\usepackage{booktabs}
\usepackage{graphicx}
\usepackage{caption}
\usepackage{subcaption}
\usepackage[backend=biber,style=numeric,sorting=none,maxbibnames=4,giveninits=true]{biblatex}
\addbibresource{references.bib}

\captionsetup{font=small,labelfont={small,sc},textfont=small}
\setcapindent{0pt}

% ---------------------------------------------------------------- notation
% Symbols follow Galvan Fraile (2025), Sections 1.2-1.3.
\newcommand{\dt}{\delta t}
\newcommand{\Isyn}{I^{\mathrm{syn}}}
\newcommand{\Csyn}{C^{\mathrm{syn}}}
\newcommand{\Iext}{I^{\mathrm{ext}}}
\newcommand{\Iint}{I^{\mathrm{int}}}
\newcommand{\amem}{\alpha}
\newcommand{\asyn}{\alpha_{r}}
\newcommand{\Acoef}{\mathcal{A}}
\newcommand{\Bcoef}{\mathcal{B}}
\newcommand{\Dcoef}{\mathcal{D}}
\newcommand{\Rcoef}{\mathcal{R}}
\newcommand{\Scoef}{\mathcal{S}}
\newcommand{\code}[1]{\texttt{\small #1}}
% classicthesis keeps a wide margin for notes; figures may use it.
\newlength{\fullwidth}
\setlength{\fullwidth}{\dimexpr\textwidth+\marginparwidth+\marginparsep\relax}
\newcommand{\widefigure}[2][\fullwidth]{%
  \makebox[\textwidth][l]{\includegraphics[width=#1]{#2}}}
% A table may use the margin too; captured in a box so \makebox stays balanced.
\newsavebox{\widetablebox}
\newenvironment{widetable}%
  {\begin{lrbox}{\widetablebox}\begin{minipage}{\fullwidth}\small\centering}%
  {\end{minipage}\end{lrbox}%
   \noindent\makebox[\textwidth][l]{\usebox{\widetablebox}}}

% Generated by export_figures.py - do not edit.
\newcommand{\nCellTypes}{201}
\newcommand{\tauMemMin}{5.5}
\newcommand{\tauMemMed}{14.3}
\newcommand{\tauMemMax}{51.8}
\newcommand{\tauSynFast}{0.99}
\newcommand{\tauSynSlow}{5.65}
\newcommand{\nBasis}{4}
\newcommand{\expmBound}{10^{-16}}
\newcommand{\eulerErrSub}{1.7\times 10^{-2}}
\newcommand{\exactErrSub}{4.2\times 10^{-7}}
\newcommand{\eulerErrSpk}{6.4\times 10^{-2}}
\newcommand{\exactErrSpk}{1.1\times 10^{-7}}
\newcommand{\eulerErrSpkHalf}{6.4\times 10^{-2}}
\newcommand{\exactErrSpkHalf}{3.6\times 10^{-3}}
\newcommand{\exactErrSpkHalfState}{1.0\times 10^{-3}}
\newcommand{\eulerErrSpkHalfState}{6.4\times 10^{-2}}
\newcommand{\constantPenalty}{3.6}
\newcommand{\gainSpk}{560\,742}
\newcommand{\gainSub}{39\,871}
\newcommand{\gainSpkHalf}{18}
\newcommand{\matchEuler}{7.1}
\newcommand{\matchExact}{99.9}
\newcommand{\matchExactHalf}{61.9}
\newcommand{\matchEulerHalf}{5.2}
\newcommand{\withinEuler}{27.9}
\newcommand{\withinEulerHalf}{26.8}
\newcommand{\withinEulerMixed}{28.1}
\newcommand{\withinExact}{100.0}
\newcommand{\withinExactHalf}{84.3}
\newcommand{\withinExactMixed}{99.7}
\newcommand{\matchEulerMixed}{7.2}
\newcommand{\matchExactMixed}{95.3}
\newcommand{\eulerErrSpkMixed}{6.4\times 10^{-2}}
\newcommand{\exactErrSpkMixed}{3.4\times 10^{-5}}
\newcommand{\rateErrEuler}{3.5}
\newcommand{\refRate}{10.7}
\newcommand{\halfMissed}{657}
\newcommand{\halfOneOff}{59}
\newcommand{\halfWithinOne}{84}
\newcommand{\closedSubsteps}{4096}
\newcommand{\closedSubstepsCoarse}{512}
\newcommand{\matchExactCoarse}{99.7}
\newcommand{\closedRefErr}{1.6\times 10^{-5}}
\newcommand{\refSpikes}{1\,723}
\newcommand{\pspUnderestimate}{7.0}
\newcommand{\pspTauM}{16.2}
\newcommand{\hybridSubEvents}{4.2\times 10^{-7}}
\newcommand{\hybridSubCurrent}{1.7\times 10^{-2}}
\newcommand{\hybridSpkEvents}{4.9\times 10^{-2}}
\newcommand{\hybridSpkCurrent}{2.6\times 10^{-2}}
\newcommand{\resetTimingError}{0.07}
\newcommand{\ascFracOfReset}{28}
\newcommand{\eulerErrSpkPct}{6}
\newcommand{\constantWidenMB}{12}
\newcommand{\constantWords}{25}
\newcommand{\ascKError}{2.0\times 10^{-3}}
\newcommand{\ascAmplitudeError}{0.6}
\newcommand{\tRefMin}{2}
\newcommand{\tRefMax}{10}
\newcommand{\precisionGain}{29}
\newcommand{\precisionBytes}{27}
\newcommand{\exactErrSubHalf}{3.6\times 10^{-3}}
\newcommand{\exactErrSubHalfState}{6.6\times 10^{-4}}
\newcommand{\eulerErrSubHalf}{1.8\times 10^{-2}}
\newcommand{\eulerErrSubHalfState}{1.7\times 10^{-2}}
\newcommand{\kernelRatio}{1.09}
\newcommand{\kernelRatioMin}{1.06}
\newcommand{\kernelRatioMax}{1.11}
\newcommand{\kernelNeurons}{230\,924}
\newcommand{\kernelBatch}{32}
\newcommand{\kernelConstantMB}{4.6}
\newcommand{\gradWorst}{1.9\times 10^{-2}}
\newcommand{\gradWorstPrevZ}{8.1\times 10^{-5}}
\newcommand{\gradBackends}{2}


\begin{document}

\title{\rmfamily\normalfont\spacedallcaps{Exact Propagators for the GLIF3 Step}}
\subtitle{\rmfamily\normalfont\spacedlowsmallcaps{%
Replacing the Euler step on the synaptic current in the V1 model}}
\author{}
\date{}
\maketitle
\vspace{-3.5em}

\begin{abstract}
\noindent
The V1 cortical model advances each GLIF3 neuron with a step of
$\dt = 1\,$ms. The synaptic and after-spike currents obey linear equations,
and the original code already used their exact one-step solutions; only the
membrane update was approximate. It used a
\emph{Euler} step, which freezes the total synaptic current at its
start-of-step value for the whole millisecond, and applies the reset and the
after-spike current of a spike one step late. With a fastest synaptic basis
of $\tau_1 = \tauSynFast\,$ms, that current rises and falls by an order of
magnitude \emph{inside} one step. Together the two approximations are the
dominant error of the simulation --- not, as one might assume, the
\code{float16} arithmetic used during training --- and at $\sim$10\,Hz the
event timing contributes the larger share. We derive the exact one-step propagator for the coupled
(synaptic, after-spike, membrane) system, following the construction of
Rotter and Diesmann~\cite{rotter1999exact}, and fold the spike events into
the same constants. Against a float64 ground truth obtained by a third,
independent method, the relative RMS error of the membrane potential falls by
a factor of $\gainSpk$, and the fraction of spikes reproduced at exactly the
right millisecond rises from $\matchEuler\,\%$ to $\matchExact\,\%$. The new
step requires no additional state variables and costs $\kernelRatio\times$
the time of the fused CUDA kernel it replaces. In the \code{float16}
configuration training actually uses, the gain is $\gainSpkHalf\times$ rather
than $\gainSpk$, and the binding constraint is no longer the integrator but
the \code{float16} storage of the propagator constants --- which, being shared
across the batch and the sequence, would cost $\constantWidenMB\,$MB in total
to widen.
\end{abstract}

\vspace{1em}
\tableofcontents
\vspace{1em}

% =========================================================================
\section{The model and its discretisation}
\label{sec:model}

\subsection{Continuous-time dynamics}

We follow the notation of~\cite{galvan2025thesis}. A GLIF3
neuron~\cite{mihalas2009generalized,teeter2018generalized} carries a membrane
potential $V(t)$ and two after-spike currents (ASCs) $I^{s}(t)$,
$s\in\{1,2\}$, obeying
%
\begin{equation}
\label{eq:glif3}
\tau_m\frac{\mathrm{d}V(t)}{\mathrm{d}t}
  = -\bigl(V(t)-E_L\bigr) + R_m I(t),
\qquad
\frac{\mathrm{d}I^{s}(t)}{\mathrm{d}t} = -k^{s}I^{s}(t),
\end{equation}
%
with $\tau_m = R_m C_m$ and $I(t) = \Iint(t) + \Iext(t)$ the total input
current, where $\Iint(t)=\sum_s I^{s}(t)$
(cf.~\cite{galvan2025thesis}, Eq.~1.8). At a spike time $t$, where
$V(t) > V_{th}$,
%
\begin{equation}
\label{eq:reset}
V(t^{+}) \leftarrow E_L,
\qquad
I^{s}(t^{+}) \leftarrow I^{s}(t) + \delta I^{s},
\end{equation}
%
with $\delta I^{s}$ the constant ASC increment
(cf.~\cite{galvan2025thesis}, Eq.~1.9).

The external current is generated by current-based synapses with
alpha-function kernels. Writing $\Csyn \equiv \dot{\Isyn} + \Isyn/\tau_s$
reduces the second-order kernel to the first-order
pair~\cite{rotter1999exact,sterratt2023principles}
%
\begin{equation}
\label{eq:alpha}
\frac{\mathrm{d}\Isyn}{\mathrm{d}t} = -\frac{\Isyn}{\tau_s} + \Csyn,
\qquad
\frac{\mathrm{d}\Csyn}{\mathrm{d}t} = -\frac{\Csyn}{\tau_s}
  + \frac{e}{\tau_s}\,W_{ji}\,z_i(t - d_{ji}),
\end{equation}
%
(cf.~\cite{galvan2025thesis}, Eq.~1.25), with $z_i$ the presynaptic spike
train and $d_{ji}$ the axonal delay.

In the V1 model~\cite{billeh2020systematic} each synapse is decomposed over
$R=\nBasis$ receptor bases with fixed time constants $\tau_r$, so every
neuron $j$ carries a pair $(\Csyn_{j,r},\Isyn_{j,r})$ per basis. The total
input current is
%
\begin{equation}
\label{eq:total}
I_j(t) = \sum_{r=1}^{R} \Isyn_{j,r}(t) + \sum_{s=1}^{2} I^{s}_j(t).
\end{equation}
%
Normalising voltages by $\Delta V = V_{th}-E_L$ so that $E_L = 0$,
$V_{th} = 1$ --- as the implementation does --- gives
$R_m = 1/g$, with $g$ the leak conductance, and
Eq.~\eqref{eq:glif3} becomes $\dot V = -V/\tau_m + I/C_m$.

\subsection{The state vector}

Collecting the $2R+3$ variables of one neuron into
%
\begin{equation}
\label{eq:statevec}
\bm{y}_j = \bigl(\Csyn_{j,1},\dots,\Csyn_{j,R},\;
                 \Isyn_{j,1},\dots,\Isyn_{j,R},\;
                 I^{1}_j,\,I^{2}_j,\; V_j\bigr)^{\!\top},
\end{equation}
%
Eqs.~\eqref{eq:glif3}--\eqref{eq:alpha} read $\dot{\bm y}_j = \bm{M}_j\bm{y}_j$
between spikes, with a constant, lower-triangular generator $\bm{M}_j$. This
is the observation on which everything below rests: \emph{between spikes the
GLIF3 neuron is a linear time-invariant system}, so its one-step transition
has a closed form and need not be approximated at all.

% =========================================================================
\section{The previous scheme, and why it fails here}
\label{sec:previous}

\subsection{What was implemented}

The implementation advanced the state over $[t, t+\dt)$ as
%
\begin{align}
\label{eq:old-csyn}
\Csyn_{j,r}[t+\dt] &= \asyn\, \Csyn_{j,r}[t]
  + \frac{e}{\tau_r}\sum_{\mathrm{src}}\sum_i
    W^{\mathrm{src}}_{ji}\, z^{\mathrm{src}}_i[t-d_{ji}],
\\
\label{eq:old-isyn}
\Isyn_{j,r}[t+\dt] &= \asyn\, \Isyn_{j,r}[t] + \dt\,\asyn\, \Csyn_{j,r}[t],
\\
\label{eq:old-asc}
I^{s}_j[t+\dt] &= \beta_s\, I^{s}_j[t] + S_j[t]\,\delta I^{s},
\\
\label{eq:old-v}
V_j[t+\dt] &= \amem V_j[t] + (1-\amem)R_m\, I_j[t] - S_j[t]\,\Delta V,
\end{align}
%
with the decay factors $\amem = e^{-\dt/\tau_m}$,
$\asyn = e^{-\dt/\tau_r}$ and $\beta_s = e^{-k^{s}\dt}$, and $S_j[t]$ the
spike indicator of~\cite{galvan2025thesis}, Eq.~1.4. Equations
\eqref{eq:old-csyn}--\eqref{eq:old-v} are Eqs.~1.27 and~1.10 of the thesis,
specialised to $E_L = 0$.

Two points deserve emphasis before we criticise anything.

First, Eqs.~\eqref{eq:old-csyn}--\eqref{eq:old-asc} are \emph{already
exact}. For $\Csyn$ and $I^{s}$ this is plain exponential decay. For $\Isyn$
it is less obvious: the term $\dt\,\asyn\Csyn_{j,r}[t]$ in
Eq.~\eqref{eq:old-isyn} looks like $\Csyn$ held constant over the step, even
though $\Csyn$ decays during it. Over one step with no arriving spike,
Eq.~\eqref{eq:alpha} gives $\Csyn(t+u) = e^{-u/\tau_r}\Csyn(t)$ for
$u\in[0,\dt]$. Substituting
%
\begin{equation}
\label{eq:isyn-exact}
\Isyn(t+u) = e^{-u/\tau_r}\bigl(\Isyn(t) + u\,\Csyn(t)\bigr)
\end{equation}
%
into Eq.~\eqref{eq:alpha} gives
$\frac{\mathrm{d}}{\mathrm{d}u}\Isyn = -\Isyn/\tau_r + e^{-u/\tau_r}\Csyn(t)
= -\Isyn/\tau_r + \Csyn(t+u)$. Eq.~\eqref{eq:isyn-exact} therefore solves the
equation exactly, time-varying $\Csyn$ included. At $u=\dt$ it is
Eq.~\eqref{eq:old-isyn}. It is exact because $\Csyn$ and $\Isyn$ share the
time constant $\tau_r$. Once that common decay is factored out,
$e^{u/\tau_r}\Csyn(t+u) = \Csyn(t)$ really is constant over the step, and
the factor $\dt$ is its exact integral. As an example, take
$\tau_r = \dt = 1\,$ms, $\Isyn(t)=0$ and $\Csyn(t)=1$. Then
Eq.~\eqref{eq:old-isyn} gives $\Isyn(t+\dt) = e^{-1} \approx 0.37$, the
exact value. A genuine forward-Euler step,
$\Isyn + \dt(\Csyn - \Isyn/\tau_r)$, would give $1$. Arriving spikes enter
$\Csyn$ at a step boundary, which is exact for events on the time grid
(Section~\ref{sec:arrival}). After every step, all synaptic and after-spike
states hold their exact values, and nothing in their update needs
improving.

Second, Eq.~\eqref{eq:old-v} is not forward Euler either: the homogeneous
part carries the exact factor $e^{-\dt/\tau_m}$ rather than
$1-\dt/\tau_m$. It is an \emph{exponential Euler} step, and it would be exact
if $I_j$ were constant across the step. The approximation is entirely in that
last clause. We refer to Eq.~\eqref{eq:old-v} throughout as the \emph{Euler}
scheme, following the thesis, and note that what it approximates is a
\emph{zero-order hold} on the current: $I_j$ is held at its start-of-step value and the factor
$1-\amem = \int_0^{\dt} e^{-(\dt-u)/\tau_m}\,\mathrm{d}u / \tau_m$
is the integral of the membrane kernel against that constant.

\subsection{When an arriving spike takes effect}
\label{sec:arrival}

In Eq.~\eqref{eq:old-csyn} the new input is added to $\Csyn$ without being
decayed, and Eq.~\eqref{eq:old-isyn} uses the old $\Csyn_{j,r}[t]$. Both
are exact if the spike arrives at the \emph{end} of the step, $t+\dt$. The
delta in Eq.~\eqref{eq:alpha} makes $\Csyn$ jump by $eW_{ji}/\tau_r$ at the
arrival time, while $\Isyn$, the integral of $\Csyn$, is continuous. A jump at
$t+\dt$ has no time left to decay within the step and has not yet
contributed to $\Isyn$. From the next step on, it decays and feeds $\Isyn$
through Eq.~\eqref{eq:isyn-exact}.

Placing the arrival at the end of the step is what makes the delays come out
right. Let a presynaptic neuron $i$ spike at grid point $t_n$. That spike is
computed from $V_i[t_n]$, and the synapse onto $j$, with delay
$d_{ji} = d\,\dt$, reads it $d-1$ steps later. The jump then lands at the end
of that step, at $t_n + d\,\dt$. This is exactly the delay in the data. Had
the jump been placed at the start of the same step, every synapse would act
one step early, and the shortest delay of $1\,$ms would become
instantaneous. Both conventions solve Eq.~\eqref{eq:alpha} exactly; they
differ only in the arrival time they assign. The afferent LGN input follows
the same rule.

The neuron's own spike is the opposite case. It happens at $t_n$, the start
of the step that follows it, so its reset and after-spike increment belong
at the start of that step (Section~\ref{sec:timing}):
%
\begin{center}
\begin{tabular}{lll}
\toprule
event & occurs at & enters the step \\
\midrule
own spike: reset and $\delta I^{s}$ & $t_n$ & at the start of $[t_n, t_{n+1})$ \\
arrival at $j$ with delay $d\,\dt$ & $t_n + d\,\dt$ & at the end of $[t_{n+d-1}, t_{n+d})$ \\
\bottomrule
\end{tabular}
\end{center}

\subsection{Why holding the current is not a small approximation}

Holding the current constant is defensible when it varies slowly
compared with $\dt$. Figure~\ref{fig:timescales} shows that in this model it
does not. The four receptor bases have
$\tau_r \in \{\tauSynFast, 1.88, 3.63, \tauSynSlow\}\,$ms against
$\dt = 1\,$ms, while the membrane constants span
$\tauMemMin$--$\tauMemMax\,$ms (median $\tauMemMed\,$ms) across the
$\nCellTypes$ cell types. The alpha kernel driven by the fastest basis peaks
at $t = \tau_1 \approx \dt$: it rises from zero to its maximum and begins to
decay \emph{within the first step after the presynaptic spike}.

\begin{figure}[t]
  \centering
  \widefigure[\fullwidth]{figures/timescales.pdf}
  \caption[Timescales versus the integration step]{%
    \textbf{The synaptic timescales collide with the step.}
    \emph{Left:} distribution of membrane time constants $\tau_m$ over the
    $\nCellTypes$ GLIF3 cell types (blue), with the four receptor-basis
    constants $\tau_r$ (orange) and the step $\dt$ (dotted).
    \emph{Right:} the alpha current $\Isyn_{j,1}(t)$ produced by one
    presynaptic spike on the fastest basis (blue), against the staircase the
    Euler step integrates in its place (orange). Because
    $\Isyn_{j,r}(t_s) = 0$ at the arrival time, Euler assigns the spike's
    own step a current of exactly zero.}
  \label{fig:timescales}
\end{figure}

The right panel of Figure~\ref{fig:timescales} makes the failure concrete.
A newly arrived spike sets $\Csyn_{j,r}$, not $\Isyn_{j,r}$; the latter is
still zero at the grid point. Euler therefore integrates the entire first
millisecond of every postsynaptic current as \emph{zero}, and then lags by one
step thereafter.

\subsection{Two further inconsistencies in the event timing}
\label{sec:timing}

The same audit exposes two off-by-one-step errors in how spike events enter
Eqs.~\eqref{eq:old-asc}--\eqref{eq:old-v}. A postsynaptic spike is emitted at
grid point $t$ from $V_j[t]$, so its consequences act \emph{at} $t$ and must
then propagate across $[t,t+\dt)$:
%
\begin{itemize}
  \item the reset in Eq.~\eqref{eq:old-v} subtracts $\Delta V$ at the
        \emph{end} of the step, giving $\amem V_j[t] - \Delta V$. A reset at
        $t$ gives $\amem\bigl(V_j[t]-\Delta V\bigr)$ instead, a discrepancy of
        $(1-\amem)\Delta V \approx \resetTimingError\,\Delta V$ at the median
        $\tau_m$;
  \item the ASC increment in Eq.~\eqref{eq:old-asc} is likewise added after
        the decay, and the freshly injected ASC does not drive $V_j$ during
        the step in which it appears.
\end{itemize}

\subsection{Event conventions against the fitted model}
\label{sec:conventions}

The GLIF3 parameters were fitted by Teeter
\emph{et al.}~\cite{teeter2018generalized}, and the fitting procedure fixes
where each spike event belongs. The fitting code\footnote{%
\url{https://github.com/AllenInstitute/AllenSDK/blob/master/allensdk/internal/model/glif/ASGLM.py}
and \code{create\_basis\_IPSP} in
\url{https://github.com/AllenInstitute/AllenSDK/blob/master/allensdk/internal/model/GLM.py}.}
regresses the sub-threshold voltage on basis currents
$e^{-k^{s}(t-t_{\mathrm{sp}})}$ that start at spike initiation $t_{\mathrm{sp}}$,
and only removes the samples inside the refractory window (the ``spike
cut'') from the regression. The fitted ASC is therefore
$\delta I^{s} e^{-k^{s}(t-t_{\mathrm{sp}})}$ from the spike onwards, and a current
present before the spike decays through the window, as in their
supplementary Eq.~6. The membrane is not modelled inside the window and
restarts at $V_{\mathrm{reset}} = E_L$ when it ends.

The exact step follows this on the grid. A spike detected at $t$ resets the
membrane and injects $\delta I^{s}$ at the start of $[t, t+\dt)$, and both
then evolve through the step (Eq.~\eqref{eq:events}). Further spikes are
blocked until $t + \lceil t_{\mathrm{ref}}/\dt\rceil\dt$. Because the true
crossing lies in $(t-\dt, t]$, that is the first grid point that can follow
the end of the refractory window. Under a hard reset the membrane is held at
$V_{\mathrm{reset}}$ up to the grid point before it. A trace of one spike through
both backends, for both resets, reproduces these values to within
$10^{-6}$ of the corresponding matrix-exponential solution.

The implementation departs from Eq.~\eqref{eq:reset} in one deliberate
respect. For training it resets by subtraction, $V \leftarrow V - \Delta V$,
rather than to $E_L$, and lets the membrane integrate its input during the
refractory period while only spiking is blocked. This keeps the overshoot
$V_j[t]-V_{th}$ and the information the input carries through the
refractory period. It is also a departure from the fit: the after-spike
currents keep driving the membrane during the refractory period, while the
fitted model does not simulate the membrane there and the hard reset holds
it. Only with no input, no overshoot and no after-spike current does the
soft reset return to $E_L$ as the hard reset does. The hard reset, which
reproduces the fitted model on the $\dt$ grid, remains available. It used to fall back to the subtractive rule when the
refractory period was a single step; it now resets to $V_{\mathrm{reset}}$ in that
case too. The current network is not affected, since its shortest
$t_{\mathrm{ref}}$ is $1.05\,$ms.

NEST's \code{glif\_psc}~\cite{gewaltig2007nest}\footnote{%
\url{https://github.com/nest/nest-simulator/blob/master/models/glif_psc.cpp},
function \code{glif\_psc::update}.} also resets at the grid point where the
crossing is detected, so the reset timing agrees. It differs in two ways,
both inherited from the Allen simulator. It adds $\delta I^{s}$ undecayed at
the \emph{end} of the refractory window, which does not match the fitted
current above. It also counts the refractory period from the detection time
as if the crossing happened exactly there, which delays the next spike by one
step whenever the fractional part of $t_{\mathrm{ref}}$ is zero or at least
one half.

% =========================================================================
\section{The exact step}
\label{sec:exact}

\subsection{Propagator}

Since $\dot{\bm y}_j = \bm{M}_j \bm{y}_j$ is linear and time invariant, the
exact transition over one step is the matrix exponential
$\bm{y}_j[t+\dt] = e^{\bm{M}_j \dt}\,\bm{y}_j[t]$. This is the
propagator-matrix construction of Rotter and
Diesmann~\cite{rotter1999exact}, used by NEST for
\code{iaf\_psc\_alpha}~\cite{gewaltig2007nest}. Only the last row is new: the
first $2R+2$ rows reproduce Eqs.~\eqref{eq:old-csyn}--\eqref{eq:old-asc}.

Integrating the membrane equation exactly rather than under a hold gives a
convolution in place of a product,
%
\begin{equation}
\label{eq:convolution}
V_j(t+\dt) = e^{-\dt/\tau_m} V_j(t)
  + \frac{1}{C_m}\int_{0}^{\dt}
      e^{-(\dt-u)/\tau_m}\, I_j(t+u)\,\mathrm{d}u,
\end{equation}
%
and because each source in Eq.~\eqref{eq:total} is a known exponential or
alpha function on $[t,t+\dt)$, every integral is a constant. Writing
$a = 1/\tau_m$ and $b_r = 1/\tau_r$, the resulting step is
%
\begin{equation}
\label{eq:exact-v}
\boxed{\;
\begin{aligned}
V_j[t+\dt] = {} & \amem V_j[t]
  + \sum_{r=1}^{R}\Bigl(\Acoef_{j,r}\,\Isyn_{j,r}[t]
                      + \Bcoef_{j,r}\,\Csyn_{j,r}[t]\Bigr) \\[0.2em]
  & + \sum_{s=1}^{2}\Dcoef^{s}_{j}\,I^{s}_j[t]
  + \bigl(\Rcoef_j + \Scoef_j\bigr)\, S_j[t]
\end{aligned}\;}
\end{equation}
%
with the coefficients
%
\begin{align}
\label{eq:coefA}
\Acoef_{j,r} &= \frac{1}{C_m}\int_0^{\dt} e^{-a(\dt-u)}e^{-b_r u}\,\mathrm{d}u
   = \frac{1}{C_m}\,\frac{e^{-b_r \dt} - e^{-a\dt}}{a - b_r},
\\
\label{eq:coefB}
\Bcoef_{j,r} &= \frac{1}{C_m}\int_0^{\dt} e^{-a(\dt-u)}\,u\,e^{-b_r u}\,\mathrm{d}u
   = \frac{1}{C_m}\left[\frac{\dt\,e^{-b_r\dt}}{a-b_r}
     - \frac{e^{-b_r\dt}-e^{-a\dt}}{(a-b_r)^{2}}\right],
\\
\label{eq:coefD}
\Dcoef^{s}_{j} &= \frac{1}{C_m}\int_0^{\dt} e^{-a(\dt-u)}e^{-k^{s}u}\,\mathrm{d}u
   = \frac{1}{C_m}\,\frac{e^{-k^{s}\dt} - e^{-a\dt}}{a - k^{s}} .
\end{align}

\subsection{Spike events}

A postsynaptic spike does two distinct things to its own step, and applying
both at the spike's grid point, then propagating them across the step,
gives, for the subtractive reset the implementation uses
(Section~\ref{sec:conventions}),

\begin{equation}
\label{eq:events}
\underbrace{\Rcoef_j = -\amem\,\Delta V}_{\text{the reset}},
\qquad
\underbrace{\Scoef_j = \sum_{s=1}^{2}\Dcoef^{s}_{j}\,\delta I^{s}}_{%
  \text{the ASC it injects}},
\quad
I^{s}_j[t+\dt] = \beta_s I^{s}_j[t]
  + \underbrace{\beta_s\,\delta I^{s}}_{\widetilde{\delta I^{s}}}\,S_j[t].
\end{equation}

$\Rcoef_j$ is the reset, which now decays across the step rather than being
subtracted at its end. $\Scoef_j$ is the drive the freshly injected after-spike
current contributes to the membrane during that same step --- a term the
Euler scheme did not have at all, because there the ASC only reached the
membrane on the following step.

Both multiply the same $S_j[t]$, so the forward pass could fold them into one
constant and save an operation. They are deliberately kept apart, because
the two training flags gate them separately. The flag \code{detach\_reset}
cuts the gradient path from the spike to the reset, and
\code{detach\_asc\_reset} cuts the path from the spike to the ASC. Folding them would make \code{detach\_reset} silently cut
part of the ASC gradient too --- and $\Scoef_j$ reaches $\ascFracOfReset\,\%$ of
$|\Rcoef_j|$ for some cell types, so that is not a rounding matter. Apart from
that one extra multiply-add, Eq.~\eqref{eq:exact-v} has the same operation
count on the spike term as Eq.~\eqref{eq:old-v}.

\subsection{Removable singularities}

Both $\Acoef_{j,r}$ and $\Dcoef^{s}_j$ carry a removable singularity at
$a = b_r$ (respectively $a = k^{s}$), and $\Bcoef_{j,r}$ a double one. This
is not a hypothetical concern: the slowest basis, $\tauSynSlow\,$ms, lies
inside the membrane range $\tauMemMin$--$\tauMemMax\,$ms, so some cell types
land on the degeneracy to within rounding. Writing $x = (a-b_r)\dt$ and
%
\begin{equation}
\label{eq:stable}
\Acoef_{j,r} = \frac{\dt\, e^{-a\dt}}{C_m}\cdot\frac{e^{x}-1}{x},
\qquad
\Bcoef_{j,r} = \frac{\dt^{2} e^{-a\dt}}{C_m}\cdot
  \frac{e^{x}(x-1)+1}{x^{2}},
\end{equation}
%
both factors are entire, and at $x = 0$ reduce to
$\Acoef = \dt\,e^{-a\dt}/C_m$ and $\Bcoef = \tfrac12 \dt^{2}e^{-a\dt}/C_m$.

Being entire is not the same as being computable, and the second factor is a
trap. Written as it stands, $e^{x}(x-1)+1$ subtracts two quantities of size
$1$ to produce one of size $x^{2}/2$, so it forfeits $2\log_{10}(1/|x|)$
significant digits --- eight of them at $x = 10^{-4}$. The band is reached in
practice: across the $\nCellTypes$ cell types the smallest $|x|$ is
$1.8\times10^{-3}$. Rewriting through
%
\begin{equation}
\label{eq:phi2}
\frac{e^{x}(x-1)+1}{x^{2}} = 1 + (x-1)\,\varphi_2(x),
\qquad
\varphi_2(x) = \frac{1}{x}\left(\frac{\mathrm{e}^{x}-1}{x} - 1\right),
\end{equation}
%
leaves a cancellation of $1$ against $1/2$ rather than against $x^{2}/2$, which
costs one order of $x$ instead of two. Below a crossover the Taylor series
$\sum_n x^{n}(n+1)/(n+2)!\;=\;\tfrac12 + \tfrac{x}{3} + \tfrac{x^{2}}{8}
+ \tfrac{x^{3}}{30} + \tfrac{x^{4}}{144} + \dots$ takes over; the crossover is
placed where the series truncation meets the recurrence's rounding. The worst
relative error over all $x$ falls from $1.8\times10^{-8}$ to
$3.6\times10^{-14}$, and $\Bcoef$ then matches
$e^{\bm{M}_j\dt}$ to $10^{-16}$ for every cell type. The first factor,
$\tfrac{e^x-1}{x} = 1 + \tfrac{x}{2} + \tfrac{x^2}{6} + \dots$, needs no such
care: \code{expm1} evaluates it to machine precision directly.

\subsection{The old scheme as a special case}

Setting
%
\begin{equation}
\label{eq:euler-special}
\begin{gathered}
\Acoef_{j,r} = \Dcoef^{s}_{j} = (1-\amem)R_m,
\qquad
\Bcoef_{j,r} = 0, \\[0.25em]
\Rcoef_j = -\Delta V,
\qquad
\Scoef_j = 0,
\qquad
\widetilde{\delta I^{s}} = \delta I^{s}
\end{gathered}
\end{equation}
%
turns Eq.~\eqref{eq:exact-v} back into Eq.~\eqref{eq:old-v} exactly. The
integration scheme is therefore a \emph{choice of constants}, not a branch in
the code: both the TensorFlow path and the fused CUDA kernel evaluate
Eq.~\eqref{eq:exact-v} unconditionally, and the historical behaviour is
recovered bit for bit by building the coefficients of
Eq.~\eqref{eq:euler-special} at setup time.

\subsection{Gradients}
\label{sec:gradients}

The model is trained by backpropagation through time, so the propagator is only
as useful as its Jacobian. Differentiating Eq.~\eqref{eq:exact-v} and the state
updates gives, for upstream gradients $\bar V, \bar I^{s}, \bar C, \bar I$ on
the four outputs,

\begin{equation}
\label{eq:jacobian}
\begin{aligned}
\frac{\partial \mathcal{L}}{\partial V_j[t]} &= \amem\,\bar V, &
\frac{\partial \mathcal{L}}{\partial \Isyn_{j,r}[t]}
  &= \asyn\,\bar I_{r} + \Acoef_{j,r}\,\bar V, \\[0.25em]
\frac{\partial \mathcal{L}}{\partial I^{s}_j[t]}
  &= \beta_s\,\bar I^{s} + \Dcoef^{s}_{j}\,\bar V, &
\frac{\partial \mathcal{L}}{\partial \Csyn_{j,r}[t]}
  &= \asyn\,\bar C_{r} + \dt\,\asyn\,\bar I_{r}
     + \Bcoef_{j,r}\,\bar V,
\end{aligned}
\end{equation}

and for the incoming spike, with the two detach flags written as indicators,

\begin{equation}
\label{eq:zgrad}
\frac{\partial \mathcal{L}}{\partial S_j[t]}
  = \mathbb{1}_{\text{attach reset}}\,\Rcoef_j\,\bar V
  + \mathbb{1}_{\text{attach ASC}}
    \Bigl(\Scoef_j\,\bar V
        + \sum_{s}\widetilde{\delta I^{s}}\,\bar I^{s}\Bigr).
\end{equation}

Every new coefficient appears in the backward pass exactly where it appears in
the forward one, so the added cost is the same handful of multiply-adds.
Equation~\eqref{eq:zgrad} is where the separation of $\Rcoef_j$ and $\Scoef_j$
earns its keep: each flag gates its own term.

The TensorFlow path is built from plain operations and is differentiated by
autodiff, but the CUDA path carries a hand-written backward kernel, so both were
checked against central finite differences of the forward step
(\code{verify\_gradients.py}). Two inputs are excluded from that check because a
finite difference across them measures a step function rather than the
transition: the emitted spikes, which pass through a Heaviside with a surrogate
gradient, and the delayed spike buffer, which feeds the event-driven recurrent
kernel where a spike is consumed as a binary event. Both are pre-existing
choices of the model, independent of the integrator.

Over the remaining, smooth inputs --- $V_j$, $\Csyn_{j,r}$, $\Isyn_{j,r}$,
$I^{s}_j$ and the incoming spike --- both backends and both schemes agree with
finite differences to a worst relative discrepancy of $\gradWorst$, which is the
float32 difference-quotient noise floor at this step size; the spike gradient of
Eq.~\eqref{eq:zgrad}, which is exactly linear and therefore free of truncation
error, agrees to $\gradWorstPrevZ$. The CUDA and TensorFlow gradients agree with
each other to float32 rounding, which the test suite asserts independently.

\subsection{Verification of the closed forms}

Equations~\eqref{eq:coefA}--\eqref{eq:coefD} were checked against
$e^{\bm{M}_j\dt}$ computed by scaling-and-squaring~\cite{higham2005scaling}
(\code{scipy.linalg.expm}) for every fifth cell type. The largest discrepancy
in any coefficient is below $\expmBound$ --- machine precision, and an
independent confirmation of the algebra, since the two routes share no code.
(Its last digits depend on the LAPACK build, so only the order is quoted.)

% =========================================================================
\section{Accuracy}
\label{sec:accuracy}

\subsection{An independent ground truth}

Measuring the exact scheme against a reference built \emph{from} the exact
propagator would be circular. We therefore construct the ground truth from the
\emph{old} method: the Euler step is first-order accurate, so
sub-stepping it at $\dt/M$ converges to the true trajectory as
$\mathcal{O}(1/M)$, and a Richardson extrapolation~\cite{richardson1911approximate}
%
\begin{equation}
\label{eq:richardson}
V^{\mathrm{ref}} = 2\,V^{\mathrm{euler}}_{2M} - V^{\mathrm{euler}}_{M},
\qquad M = 128,
\end{equation}
%
cancels the leading error term, leaving a reference accurate to
$\sim\!10^{-7}$ that never touches the closed form under test. The reference
takes only Euler's \emph{integration} from the old method: it applies the
spike events at the start of each step, as the exact scheme does
(Section~\ref{sec:conventions}), so it measures only the integration error.
The event convention is part of the model's definition, fixed by the fit,
not an approximation a convergent method could arbitrate. All errors
below are relative RMS errors in $V_j$, in units of $\Delta V$, over
$\nCellTypes$ cell types driven by a fixed Poisson input.

Figure~\ref{fig:convergence} shows the sub-stepped Euler step walking towards the
propagator as $1/M$ and never overtaking it. The propagator is already at the
reference's own accuracy floor at $M = 1$.

\begin{figure}[t]
  \centering
  \widefigure[0.78\fullwidth]{figures/convergence.pdf}
  \caption[Convergence of the sub-stepped Euler step]{%
    \textbf{Euler converges to the propagator.}
    Relative RMS error in $V_j$ against the Richardson reference of
    Eq.~\eqref{eq:richardson}, as the Euler step is sub-stepped at
    $\dt/M$ (orange circles). The square marks the scheme as implemented,
    which additionally carries the end-of-step event timing of
    Section~\ref{sec:previous}. The exact propagator (blue) sits at the
    reference's own floor for $M=1$.}
  \label{fig:convergence}
\end{figure}

\subsection{Error in the two regimes, at both precisions}

Table~\ref{tab:errors} and Figure~\ref{fig:errors} separate two effects. The
\emph{subthreshold} regime has no postsynaptic spikes, so the error is purely
the membrane integration of the synaptic current; the \emph{spiking} regime
adds the reset and the ASCs, whose event timing the exact scheme also
corrects. Each is evaluated in float64, which isolates the method, and with
the state stored in \code{float16} and the arithmetic performed in
\code{float32} --- exactly what the production CUDA kernel
does~\cite{micikevicius2018mixed}.

\begin{table}[t]
  \caption[Relative error of the two schemes]{%
    Relative RMS error in $V_j$ against the Richardson reference. ``state''
    narrows only the four recurrent state blocks; ``+ constants'' also narrows
    the propagator coefficients, which is what the kernel really does under a
    mixed-precision policy. The float64 entries for the exact scheme are
    bounded by the reference, not by the scheme, so the improvement factor in
    that column is a lower bound.}
  \label{tab:errors}
  \begin{widetable}
  \begin{tabular}{llrrrr}
    \toprule
    & & & \multicolumn{2}{c}{\code{float16}} & \\
    \cmidrule(lr){4-5}
    Regime & Scheme & float64 & state & + constants & Improvement \\
    \midrule
    Subthreshold      & Euler            & $\eulerErrSub$  & $\eulerErrSubHalfState$ & $\eulerErrSubHalf$  & \\
                      & exact propagator & $\exactErrSub$  & $\exactErrSubHalfState$ & $\exactErrSubHalf$  & $\gainSub\times$ \\[0.3em]
    Spiking ($\refRate\,$Hz)
                      & Euler            & $\eulerErrSpk$  & $\eulerErrSpkHalfState$ & $\eulerErrSpkHalf$  & \\
                      & exact propagator & $\exactErrSpk$  & $\exactErrSpkHalfState$ & $\exactErrSpkHalf$  & $\gainSpk\times$ \\
    \bottomrule
  \end{tabular}
  \end{widetable}
\end{table}

\begin{figure}[t]
  \centering
  \widefigure[0.84\fullwidth]{figures/errors.pdf}
  \caption[Method error versus precision error]{%
    \textbf{The method error dominated the precision error.}
    Under Euler, moving from float64 to \code{float16} storage changes
    nothing: the discretisation error is some forty times larger than
    anything the number format contributes. Only once the integrator is exact
    does \code{float16} storage become the accuracy floor.}
  \label{fig:errors}
\end{figure}

\textbf{The float64 columns are the substantive result. Euler's error,
$\bm{\eulerErrSpk}$ of $\bm{\Delta V}$, is over $\bm{\eulerErrSpkPct\,\%}$ of the distance from rest to
threshold --- a systematically wrong trajectory, not a rounding artefact.} The
three columns read together settle the question of precision, and reverse the
expected answer twice. \textbf{Mixed precision was never the \emph{integrator's}
bottleneck --- under Euler, narrowing the state changes nothing, because the
discretisation error is some forty times larger.} But once the integrator is
exact, precision becomes the floor, and it is not the state that dominates it:
narrowing the propagator \emph{constants} costs a further
$\constantPenalty\times$ on top of narrowing the state. Section~\ref{sec:floor}
returns to this, because the constants are by far the cheapest thing here to
widen.

Euler's float64 error has two sources: the current held over the step
(Section~\ref{sec:previous}), and the spike events applied at the end of the
step (Section~\ref{sec:timing}). They live in disjoint constants, the
current in $\Acoef$, $\Bcoef$, $\Dcoef$ and the events in $\Rcoef$, $\Scoef$
and the ASC increment, so each can be fixed on its own:
%
\begin{center}
\begin{tabular}{lcc}
\toprule
remaining approximation & subthreshold & spiking \\
\midrule
both (Euler)                 & $\eulerErrSub$      & $\eulerErrSpk$ \\
event timing only            & $\hybridSubEvents$  & $\hybridSpkEvents$ \\
current hold only            & $\hybridSubCurrent$ & $\hybridSpkCurrent$ \\
neither (exact)              & $\exactErrSub$      & $\exactErrSpk$ \\
\bottomrule
\end{tabular}
\end{center}
%
Without spikes the current hold is the entire error. At the imposed
$\sim$10\,Hz output rate, the event timing is the larger of the two, and
neither fix alone gets within two orders of magnitude of the exact step. The
split depends on the rate: the event error scales with the number of spikes,
the hold error with the synaptic drive.

\subsection{A single postsynaptic potential}

Figure~\ref{fig:psp} isolates the mechanism on one neuron
($\tau_m = \pspTauM\,$ms) receiving a single presynaptic spike. Euler
cannot see the current rise inside the step in which it happens, so the
postsynaptic potential starts a millisecond late and underestimates the peak
by $\pspUnderestimate\,\%$. The exact propagator is indistinguishable from
the ground truth.

\begin{figure}[t]
  \centering
  \widefigure[\fullwidth]{figures/psp.pdf}
  \caption[Postsynaptic potential]{%
    \textbf{One presynaptic spike, one neuron.}
    Membrane response to a single spike on the fastest receptor basis at
    $t = 2\,$ms. The grey trace is the sub-stepped float64 ground truth; the
    exact propagator (blue) lies on top of it. \emph{Right:} the first steps,
    where Euler (orange) integrates the arrival step as zero.}
  \label{fig:psp}
\end{figure}

\subsection{Spike-level fidelity}

Voltage error matters only insofar as it moves spikes. We therefore close the
loop: neurons integrate, cross threshold, reset, adapt and block during their
refractory period, driven by a fixed Poisson input at $\refRate\,$Hz
($\refSpikes$ reference spikes over $800\,$ms). The metric is deliberately
strict --- the fraction of ground-truth spikes reproduced in \emph{exactly}
the same millisecond --- because a closed loop amplifies any voltage
discrepancy into a different spike time, after which the trajectories
separate.

\begin{figure}[t]
  \centering
  \widefigure[\fullwidth]{figures/spikes.pdf}
  \caption[Spike-level agreement with the ground truth]{%
    \textbf{Spikes land in the right millisecond.}
    \emph{Left:} fraction of ground-truth spikes reproduced at exactly the
    right millisecond. \emph{Right:} rasters for the fourteen most active
    neurons over a $160\,$ms window, with the exact propagator in float64
    (top), the exact propagator in the production \code{float16}
    configuration (middle) and Euler in float64 (bottom); a scheme is correct
    where its tick aligns with the grey ground-truth bar. The ground truth here is the \emph{old}
    method sub-stepped at $\dt/\closedSubsteps$ in float64, so it favours
    neither scheme. It applies the spike events at the start of each step,
    as the exact scheme does, so it measures only the integration error. Its own voltage error is about $\closedRefErr$ of
    $\Delta V$, so the exact scheme's float64 agreement is a lower bound.}
  \label{fig:spikes}
\end{figure}

The result (Figure~\ref{fig:spikes}) is that Euler reproduces
$\matchEuler\,\%$ of the ground-truth spikes and the exact propagator
$\matchExact\,\%$. In the \code{float16} configuration training uses --- state
\emph{and} constants --- the exact scheme reproduces $\matchExactHalf\,\%$.
Richardson extrapolation cannot supply this ground truth: once neurons spike
on their own, runs at $M$ and $2M$ sub-steps can diverge by a whole spike,
and a raster cannot be extrapolated. It is therefore plain sub-stepped Euler
at $M = \closedSubsteps$, whose voltage error, first order in $1/M$, is about
$\closedRefErr$ of $\Delta V$. That is small next to Euler's error but not
next to the exact scheme's, so part of the exact scheme's float64 mismatch is
the reference's own: against a reference with $M = \closedSubstepsCoarse$, the
same run agrees on only $\matchExactCoarse\,\%$ of spikes. The float64
agreement of the exact scheme is therefore a lower bound.
The $\matchExactHalf\,\%$ is a strict count, and most of the remaining
spikes are not far off: of the $\halfMissed$ ground-truth spikes the
\code{float16} run misses, $\halfOneOff\,\%$ have a \code{float16} spike exactly
one millisecond away, so $\halfWithinOne\,\%$ of all ground-truth spikes are
reproduced to within $\pm 1\,$ms. That remaining gap is storage precision, not
the integrator, and
Section~\ref{sec:floor} shows most of it is recoverable for a fixed
$\constantWidenMB\,$MB.

It is worth stating what this does \emph{not} claim. Euler is stable and
its mean firing rate is only $\rateErrEuler\,\%$ off, so a network trained
under it is a perfectly self-consistent dynamical system. It is simply not
the GLIF3 system its parameters describe. Any quantity that depends on spike
timing --- synchrony, phase, latency, temporal codes, or the surrogate
gradients that flow through the threshold
crossing~\cite{bellec2020solution} --- was being computed for the wrong model.

% =========================================================================
\section{Cost}
\label{sec:cost}

\subsection{Arithmetic and memory}

Equation~\eqref{eq:exact-v} replaces one scaled sum by a weighted sum: per
neuron per step, $R$ additional fused multiply-adds, and per-basis coefficient
arrays where there was previously one scalar $(1-\amem)R_m$. The coefficients
are shared across the batch and so remain resident in L2. Critically,
$\Csyn_{j,r}$ was \emph{already} read by the kernel, which has to update it
anyway --- so the exact step introduces \textbf{no new state variables}, and
therefore no additional memory on the backpropagation-through-time tape,
which is what constrains this model in practice. The only new memory is
$\kernelConstantMB\,$MB of constants at $\kernelNeurons$ neurons in
\code{float16}, independent of batch size and sequence length.

\subsection{Measured kernel time}

Table~\ref{tab:bench} compares the fused forward and backward CUDA kernels
before and after the change. Each timing chains $64$ state transitions inside
a single \code{tf.function}, so it measures per-timestep kernel cost rather
than dispatch overhead.

% Generated by export_figures.py - do not edit.
\begin{table}[t]
  \centering
  \caption[Kernel timings]{%
    Per-timestep cost of the fused GLIF state kernels, in
    microseconds, on an NVIDIA RTX PRO 6000. Each timing chains
    64 state transitions inside one \code{tf.function}.
    \code{fp16} is the configuration used for training.}
  \label{tab:bench}
  \begin{tabular}{rrlrrrrr}
    \toprule
    & & & \multicolumn{2}{c}{forward} & \multicolumn{2}{c}{backward} & \\
    \cmidrule(lr){4-5}\cmidrule(lr){6-7}
    Neurons & Batch & dtype & Euler & exact & Euler & exact & Ratio \\
    \midrule
    51\,978 & 8 & \code{fp16} & 15.3 & 16.6 & 20.0 & 20.9 & 1.06$\times$ \\
    51\,978 & 8 & \code{fp32} & 26.4 & 33.1 & 32.7 & 36.7 & 1.18$\times$ \\
    51\,978 & 32 & \code{fp16} & 50.5 & 55.7 & 66.2 & 71.2 & 1.09$\times$ \\
    51\,978 & 32 & \code{fp32} & 116.9 & 135.0 & 141.7 & 150.8 & 1.11$\times$ \\
    230\,924 & 8 & \code{fp16} & 57.4 & 64.1 & 73.8 & 81.2 & 1.11$\times$ \\
    230\,924 & 8 & \code{fp32} & 136.8 & 152.2 & 156.2 & 165.6 & 1.08$\times$ \\
    230\,924 & 32 & \code{fp16} & 293.4 & 315.4 & 345.5 & 379.5 & 1.09$\times$ \\
    230\,924 & 32 & \code{fp32} & 580.2 & 647.1 & 658.9 & 690.8 & 1.08$\times$ \\
    \bottomrule
  \end{tabular}
\end{table}


In \code{float16} the fused state kernel costs $\kernelRatio\times$ its
previous time (range $\kernelRatioMin$--$\kernelRatioMax\times$). This is the
cost of that one kernel, not of a training step: the state transition is one
of several kernels per timestep and the recurrent CSR products dominate. An
end-to-end training measurement has not been performed.

The overhead is not irreducible. The inner loop reads four separate
$[N \times R]$ constant arrays at the same index; interleaving them into a
single array of four values per $(j,r)$ would turn four strided loads into one
\code{float4} fetch. That is a worthwhile optimisation, not a prerequisite.

% =========================================================================
\section{Consequences and open questions}
\label{sec:discussion}

\subsection{Where the accuracy floor now lies}
\label{sec:floor}

With the method error removed, \code{float16} \emph{storage} sets the floor.
Figure~\ref{fig:precision} attributes it, block by block, over the five things
the kernel stores: the four recurrent state blocks and the propagator
constants.

\begin{figure}[t]
  \centering
  \widefigure[0.92\fullwidth]{figures/precision.pdf}
  \caption[Cost of float16 storage per stored block]{%
    \textbf{What each stored block costs in \code{float16}.}
    Relative RMS error in $V_j$ for the exact scheme, as each block is moved to
    \code{float16}. The four state blocks are charged per neuron per sample;
    the constants are shared across the batch and the sequence, so they add
    nothing to that column. The dashed line is Euler with everything narrowed.}
  \label{fig:precision}
\end{figure}

Two results here are worth separating, because they point the same way.

Among the \emph{state} blocks the cost is counter-intuitive: $\Isyn$ and
$\Csyn$ --- eight of the eleven words per neuron --- cost almost nothing in
\code{float16}, while $I^{s}$, two words, costs more than both together. The
after-spike currents carry large amplitudes and decay over hundreds of
milliseconds, so their \code{float16} ulp is coarse and the error accumulates
for a long time before it decays. Widening $V_j$ and $I^{s}_j$ to
\code{float32} buys a $\precisionGain$-fold gain for $\precisionBytes\,\%$ more
state on this block --- a far smaller fraction of the full
backpropagation-through-time tape, which is dominated by the spike ring buffer.

The \emph{constants} are the better bargain, and were the easier thing to
overlook. Narrowing them costs $\constantPenalty\times$ on its own, more than
every state block put together, and it is the single largest term in the
production number of Table~\ref{tab:errors}. But unlike the state they do not
scale: there are $\constantWords$ words per neuron, shared across every sample
in the batch and every step of the sequence, so widening all of them to
\code{float32} at $\kernelNeurons$ neurons costs a fixed
$\constantWidenMB\,$MB. The obstacle is not memory but plumbing --- the fused
operator takes one dtype for every tensor it is given, so mixed state and
constant widths need a second type parameter on the op.

We quantify both trades here and take neither: the integrator fix is
unconditional, while the precision choices are deliberate decisions with their
own costs, and the constants one needs a kernel signature change. On the
evidence above the constants should be widened first.

\subsection{Two defects the rewrite exposed}
\label{sec:defects}

Reworking the setup code turned up two bugs that predate the propagator and
have nothing to do with it. Both are fixed here.

\paragraph{The ASC decay went through \code{float16}.}
$\beta_s = e^{-k^{s}\dt}$ was not computed from $k^{s}$ directly. The rate was
cast to the compute dtype, pushed through an inverse sigmoid, stored, and
pushed back through a sigmoid on read --- a parameterisation left over from
when $k$ was to be trained. Under a mixed-precision policy that round trip
happens in \code{float16} and returns $k^{s}$ with a relative error of
$\ascKError$. The error in $\beta_s$ looks negligible, at
$3.6\times10^{-5}$, but it is an error in a \emph{decay rate}: over three ASC
time constants it compounds to $\ascAmplitudeError\,\%$ of the after-spike
current's amplitude, and the slowest ASC has $\tau = 333\,$ms, so that is a
sustained bias on the slowest adaptation in the model. It is now evaluated in
float64 and cast once.

\paragraph{A neuron-aligned constant was not following the layout.}
The model can renumber neurons into a spatially local (Morton) order at
runtime while keeping checkpoints in canonical order, and
\code{\_translate\_neuron\_layout} permutes the neuron-aligned arrays when it
does. It worked from a hand-written list of attribute names, and
\code{t\_ref\_steps} was not on it --- so under a non-identity layout every
neuron got somebody else's refractory period, drawn from a range of
$\tRefMin$ to $\tRefMax$ steps. The read is live in both backends.

The second is the more instructive failure: a list of names that has to be
kept in step with a set of attributes will eventually fall out of step. The
edge-aligned translation next to it does not have this problem, because it
discovers its targets. The per-neuron constants are now built by a helper that
binds and registers in one call, so the registry cannot disagree with what
exists, and a test asserts that every constant the step reads is in it and
survives a permutation round trip.

\subsection{Reproducing earlier work}

The coefficients are derived constants, not trainable parameters, and are
deliberately excluded from the checkpoint graph, so checkpoints written before
this change still restore and are interchangeable between the two schemes.
What is \emph{not} interchangeable is the science: a network whose weights
were trained under Euler was fitted to a different dynamical system, and
its firing rates will shift when evaluated under the exact step. Earlier
results should be reproduced with the Euler coefficients of
Eq.~\eqref{eq:euler-special}, and new models retrained.

\subsection{What remains on the grid}

Spikes are still emitted and delivered only at millisecond boundaries.
Morrison \emph{et al.}~\cite{morrison2007exact} show how to carry continuous
spike times through the same propagator formalism, and
Hanuschkin \emph{et al.}~\cite{hanuschkin2010general} give a general
time-driven implementation; either would remove the remaining sub-millisecond
quantisation. Its largest visible effect is on the after-spike currents.
A spike is detected at the first grid point after the true crossing, but
$\delta I^{s}$ is injected there at full amplitude. The fitted current has
already decayed from spike initiation, so for the fastest ASC
($k^{s} = 0.3\,\mathrm{ms}^{-1}$) its value at that grid point averages only
$(1-e^{-k^{s}\dt})/(k^{s}\dt) \approx 86\,\%$ of $\delta I^{s}$. Carrying continuous spike times
is a larger change, orthogonal to this one, and it becomes the natural next target now
that the subthreshold integration is exact.

\subsection{Summary}

\begin{table}[h]
  \caption[Summary]{Summary of the change.}
  \label{tab:summary}
  \begin{widetable}
  \begin{tabular}{lll}
    \toprule
     & Euler & Exact propagator \\
    \midrule
    Membrane error (float64)        & $\eulerErrSpk$      & $\exactErrSpk$ \\
    Membrane error (\code{float16}) & $\eulerErrSpkHalf$  & $\exactErrSpkHalf$ \\
    Spikes at the right ms (float64)$^{b}$ & $\matchEuler\,\%$ ($\withinEuler\,\%$) & $\matchExact\,\%$ ($\withinExact\,\%$) \\
    Spikes at the right ms (\code{float16})$^{b}$ & $\matchEulerHalf\,\%$ ($\withinEulerHalf\,\%$) & $\matchExactHalf\,\%$ ($\withinExactHalf\,\%$) \\
    Membrane error (mixed)$^{a}$    & $\eulerErrSpkMixed$ & $\exactErrSpkMixed$ \\
    Spikes at the right ms (mixed)$^{a,b}$ & $\matchEulerMixed\,\%$ ($\withinEulerMixed\,\%$) & $\matchExactMixed\,\%$ ($\withinExactMixed\,\%$) \\
    Additional state variables      & ---                & none \\
    Fused kernel time (\code{fp16}) & $1.00\times$       & $\kernelRatio\times$ \\
    \bottomrule
  \end{tabular}
  \par\smallskip\footnotesize
  $^{a}$\,\code{float16} only for $\Isyn$ and $\Csyn$; $V_j$, $I^{s}_j$ and the
  propagator constants in float32, arithmetic in float32. Not available in the
  fused operator today, which takes one dtype for every tensor, so it is
  measured with the same per-block rounding emulation as Table~\ref{tab:errors};
  in the all-\code{float16} case that emulation reproduces the production CUDA
  kernel to $2\times10^{-7}$.\\
  $^{b}$\,In brackets, the fraction reproduced to within $\pm 1\,$ms.
  \end{widetable}
\end{table}

\vspace{1em}
\noindent
\begin{sloppypar}
The implementation lives in \code{glif\_propagators.py} under
\code{v1\_model\_utils}. The scheme is selected by the
\code{integration\_scheme} argument of \code{create\_model}, set to
\code{"exact"} or \code{"euler"}, or by the matching command-line flag. The
accompanying notebook reproduces every figure and every number in this report.
\end{sloppypar}

\printbibliography[heading=bibnumbered,title={References}]

\end{document}
